
\documentclass[12pt]{article}
\usepackage{amsmath}
\usepackage{amsfonts}
\usepackage{amssymb}
\usepackage{geometry}
\geometry{a4paper, margin=0.8in}
\pagenumbering{gobble}

\title{Homework solutions for 02YELMA \\ Chapter: Magnetostatics}
\begin{document}
\date{}
\maketitle
\vspace{-0.8in}

\begin{enumerate}

\item \textbf{Magnetic field at the center of a square loop}

Using the magnetic field of a finite straight wire:

\[
B=\frac{\mu_0 I}{4\pi R}\left(\sin\theta_1+\sin\theta_2\right)
\]

At the center of the square:

\[
\theta_1=\theta_2=45^\circ
\]

so for one side,

\[
B_{\text{one side}}
=
\frac{\mu_0 I}{4\pi R}
\left(
\sin45^\circ+\sin45^\circ
\right)
\]

\[
=
\frac{\mu_0 I}{4\pi R}
\left(
\frac{\sqrt2}{2}+\frac{\sqrt2}{2}
\right)
=
\frac{\mu_0 I\sqrt2}{4\pi R}
\]

There are 4 identical sides, therefore

\[
B=
4B_{\text{one side}}
=
\frac{\mu_0 I\sqrt2}{\pi R}
\]

Hence,

\[
\boxed{
B=\frac{\mu_0 I\sqrt2}{\pi R}
}
\]

The direction is perpendicular to the plane of the loop, determined by the right-hand rule.

\vspace{0.5cm}

\item \textbf{Long cylindrical wire with current on the outer surface}

Using Ampère’s law,

\[
\oint \vec B\cdot d\vec l=\mu_0 I_{\text{enc}}
\]

Take a circular Amperian loop of radius \(r\).

\textbf{Inside the wire (\(r<a\)):}

Since current exists only on the surface,

\[
I_{\text{enc}}=0
\]

therefore

\[
B(2\pi r)=0
\]

so

\[
\boxed{B=0}
\]

\textbf{Outside the wire (\(r>a\)):}

Now the enclosed current is the total current \(I\):

\[
B(2\pi r)=\mu_0 I
\]

thus

\[
\boxed{
B=\frac{\mu_0 I}{2\pi r}
}
\]

The direction is azimuthal (\(\hat{\phi}\)) by the right-hand rule.

\vspace{0.5cm}

\item \textbf{Circular current loop}

\begin{itemize}

\item[(a)] \textbf{Magnetic dipole moment}

The magnetic dipole moment is

\[
\vec m=I A \hat n
\]

where

\[
A=\pi R^2
\]

Since the current is counterclockwise as viewed from the positive \(z\)-axis, the direction is \(+\hat z\).

Thus,

\[
\boxed{
\vec m=I\pi R^2\hat z
}
\]

\item[(b)] \textbf{Approximate magnetic field far from the origin}

Using the dipole field formula,

\[
\vec{B}_{dip}(r,\theta)
=
\frac{\mu_0 m}{4\pi r^3}
\left[
2\cos\theta\,\hat r
+
\sin\theta\,\hat\theta
\right]
\]

Substituting

\[
m=I\pi R^2
\]

gives

\[
\boxed{
\vec B
=
\frac{\mu_0 I R^2}{4r^3}
\left[
2\cos\theta\,\hat r
+
\sin\theta\,\hat\theta
\right]
}
\qquad (r\gg R)
\]

\item[(c)] \textbf{Consistency with the exact result on the \(z\)-axis}

The exact magnetic field on the axis is

\[
B_z
=
\frac{\mu_0 I R^2}
{2(R^2+z^2)^{3/2}}
\]

For \(z\gg R\),

\[
R^2+z^2\approx z^2
\]

so

\[
B_z
\approx
\frac{\mu_0 I R^2}{2z^3}
\]

From the dipole formula on the axis,

\[
\theta=0,\quad r=z
\]

thus

\[
B=
\frac{\mu_0 m}{4\pi z^3}(2)
\]

and using

\[
m=I\pi R^2
\]

we obtain

\[
B=
\frac{\mu_0 I R^2}{2z^3}
\]

which is exactly the same result. Therefore,

\[
\boxed{\text{The dipole approximation is consistent with the exact field}}
\]

\end{itemize}

\vspace{0.5cm}

\item \textbf{Orbital magnetic dipole moment in Rutherford atom}

The equivalent current is

\[
I=\frac{\text{charge}}{\text{period}}
\]

The period is

\[
T=\frac{2\pi R}{v}
\]

thus

\[
I=\frac{e}{T}
=
\frac{ev}{2\pi R}
\]

The magnetic dipole moment is

\[
m=IA
=
\frac{ev}{2\pi R}\cdot \pi R^2
\]

which gives

\[
m=\frac{evR}{2}
\]

Since the electron has negative charge, the magnetic moment points opposite to the direction of current.

Counterclockwise motion of the electron means conventional current is clockwise, so the dipole moment points in the negative \(z\)-direction.

Hence,

\[
\boxed{
\vec m=
-\frac{evR}{2}\hat z
}
\]

\vspace{0.5cm}

\item \textbf{Uniformly magnetized sphere}

Given

\[
\vec M=M\hat z
\]

\begin{itemize}

\item[(a)] \textbf{Bound currents}

Volume bound current:

\[
\vec J_b=\nabla\times\vec M
\]

Since \(M\) is constant,

\[
\boxed{
\vec J_b=0
}
\]

Surface bound current:

\[
\vec K_b=\vec M\times\hat n
\]

For a sphere,

\[
\hat n=\hat r
\]

and using

\[
\hat z=\cos\theta\,\hat r-\sin\theta\,\hat\theta
\]

we get

\[
\vec K_b
=
M\hat z\times\hat r
=
M\sin\theta\,\hat\phi
\]

Therefore,

\[
\boxed{
\vec K_b=M\sin\theta\,\hat\phi
}
\]

\item[(b)] \textbf{Surface velocity of rotating shell}

Distance from the axis is

\[
R\sin\theta
\]

so

\[
v=\omega R\sin\theta
\]

and the direction is \(\hat\phi\).

Hence,

\[
\boxed{
\vec v=\omega R\sin\theta\,\hat\phi
}
\]

\item[(c)] \textbf{Surface current density}

Using

\[
\vec K=\sigma\vec v
\]

we obtain

\[
\boxed{
\vec K
=
\sigma\omega R\sin\theta\,\hat\phi
}
\]

\item[(d)] \textbf{Comparison of \(\vec K_b\) and \(\vec K\)}

We have

\[
K_b=M\sin\theta
\]

and

\[
K=\sigma\omega R\sin\theta
\]

These are identical if

\[
\boxed{
M=\sigma\omega R
}
\]

Thus, a uniformly magnetized sphere behaves like a rotating charged spherical shell.

\item[(e)] \textbf{Magnetic dipole moment}

For a uniformly magnetized sphere,

\[
\vec m
=
\vec M \times (\text{volume})
\]

that is,

\[
\vec m
=
\vec M
\left(
\frac{4}{3}\pi R^3
\right)
\]

Therefore,

\[
\boxed{
\vec m=
\frac{4}{3}\pi R^3 M\,\hat z
}
\]

\end{itemize}

\vspace{0.5cm}

\item \textbf{Electron as a rotating spherical shell}

From the previous result,

\[
m=\frac{4}{3}\pi R^3 M
\]
and
\[M=\sigma\omega R
\]
so
\[m=\frac{4}{3}\pi R^3 \sigma\omega R
=\frac{4}{3}\pi R^4 \sigma\omega
\]
Thus we can write angular velocity as
\[\omega=\frac{3m}{4\pi R^4 \sigma}
\]
Again from the previous result, we know that the surface velocity of rotating shell is
\[v=\omega R\sin\theta
\]
The maximum speed is at the equator where \(\sin\theta=1\), so
\[v_{\text{max}}=\omega R
\]
Substituting \(\omega\) gives
\[v_{\text{max}}=\frac{3m}{4\pi R^3 \sigma}
\]
Also, she surface charge density is
\[\sigma=\frac{e}{4\pi R^2}
\]
so
\[v_{\text{max}}=\frac{3m}{4\pi R^3 \cdot \frac{e}{4\pi R^2}}
=\frac{3m}{eR}
\]
From the previous homework, we calculated the classical radius of electron as
\[R=1.4 \times 10^{-15}\,\text{m}
\]
From experiment, the magnetic dipole moment of electron is
\[m=9.28 \times 10^{-24}\,\text{A}\cdot\text{m}^2
\]
and the elementary charge is
\[e=1.60 \times 10^{-19}\,\text{C}
\]
Substituting these values gives
\[v_{\text{max}}=\frac{3 \times 9.28 \times 10^{-24}}{1.60 \times 10^{-19} \times 1.4 \times 10^{-15}}
\approx 1.24 \times 10^{11}\,\text{m/s}
\]
which is more than 1000 times the speed of light. This is clearly contradict the principles of special relativity, indicating that the classical model of electron as a rotating spherical shell is not valid.
\end{enumerate}

\end{document}
